\scanpage{123}
% archive sheet number: 122

restriction dudit faisceau au \deleted{corps résiduel} point générique de $D_i$, alors \unsure{\deleted{o}\handwritten{b}|\deleted{o}\handwritten{o}}n peut identifier entre eux tous les groupes $\underline{Z}_{\handwritten{\unsure{\ell|l}}}(1)(\handwritten{\bar}{x})$ envisagés à un même groupe $\underline{\mathbb{Z}}_{\handwritten{\ell}}(1)(D_i)$ (qu'on peut interpréter comme groupe des sections de $\underline{Z}_{\handwritten{\ell}}(1)$ sur $D_i$ ou sur le point générique de $D_i$), et on trouve une classe \inserted{de conjugaison} bien déterminée d\deleted{\unsure{r|e}}'homomorphismes $\underline{Z}_{\handwritten{\ell}}(1)(D_i)\handwritten{\unsure{\mathbin{-}|\longrightarrow}}G$. (On trouve un résultat analogue pour les homomorphismes de produits de groupes $\underline{Z}_{\handwritten{\ell}}(1)$ envisagés dans le n\textsuperscript{o}3). Bien entendu, si $\underline{Z}_{\handwritten{\ell}}(1)$ est trivial sur l'ouvert\inserted{$X_{\handwritten{\ell}}$} des points de $X$ \deleted{Lorsque l'on a un $f:X\unsure{\mathbin{-}|\longrightarrow}Y$ comme au n\textsuperscript{o}\unsure{4|2}} de car. $\handwritten{\ne\ell}$, et si ce dernier est connexe, on peut identifier les groupes $\underline{Z}_{\handwritten{\ell}}(1)(D_i)$ au même groupe \inserted{\handwritten{\unsure{à|2}}} $\underline{Z}_{\handwritten{\ell}}(1)(X_{\handwritten{\ell}})$, \deleted{et on} soit $\underline{Z}_{\handwritten{\ell}}(1)$, et on trouve \deleted{un homomorphisme cano} pour chaque $i$ une\deleted{\unsure{homomo}} classe de conjugaison d'homomorphi\unsure{\deleted{s}m|\deleted{m}s|\deleted{s}n|\deleted{n}s}ens $\underline{Z}_{\handwritten{\ell}}(1)\handwritten{\longrightarrow}G$\handwritten{.} En général, celles-ci \deleted{sont} peuvent être distinctes pour des $i$ distincts, bien entendu. Mais relativement à un\deleted{e} morphisme $f:X\handwritten{\longrightarrow}Y$ comme dans les n\textsuperscript{o}2,\handwritten{4}, la formation des classes\inserted{de conjugaison} \deleted{précédentes} d'homomorphismes\inserted{ci-dessus} a un caractère évident de fonctorialité (qu'il y aurait lieu d'énoncer de façon précise, dans le cadre général du n\textsuperscript{o} 4).

\figure{\handwritten{\begin{tikzpicture}[x=0.65cm,y=0.65cm]
\draw (0,11.5) .. controls (0.2,10.8) and (0,9.6) .. (0.3,9) .. controls (0.4,6) and (0.2,3) .. (0.6,0);
\draw (0.7,-0.05) .. controls (1.8,0.1) and (3,0) .. (5,0) -- (9,0) .. controls (9.4,0) and (9.2,0.4) .. (9.7,0.35) .. controls (11.5,0.1) and (13,0.2) .. (15,0.5);
\end{tikzpicture}}}

Revenant alors aux conditions du présent n\textsuperscript{o}5, on peut préciser de la façon suivante l'assertion faite sur les \deleted{grou} classes de conjugaison de sous-groupes d'inertie en des assertions sur les classes de conjugaisons d'\unsure{\deleted{b}h|\deleted{h}b}\uline{omomorphism}es \uline{de modules de Tate}: supposons que $Y_{\handwritten{\ell}}$ soit connexe et que $Z_{\handwritten{\ell}}(1)$ soit trivia\unsure{\deleted{o}l|\deleted{l}o|\deleted{b}l|\deleted{l}b} sur $Y_{\handwritten{\ell}}$, \deleted{\unsure{d'où|d'on} un isomorphisme désignant} de sorte qu'on peut identifier $Z_{\handwritten{\ell}}(1)(Y_{\handwritten{\ell}})$ et les divers $\underline{Z}_{\handwritten{\ell}}(1)(D_i)$ à un même module \handwritten{(}libre de rang 1\handwritten{)}$\underline{Z}_{\handwritten{\ell}}(1)$ sur $\underline{Z}_{\handwritten{\ell}}$. Considérons alors les classes de conjugaison d'homomorphismes composés
\diagram{\begin{tikzcd}[ampersand replacement=\&,column sep=1.4cm]
\underline{Z}_{\handwritten{\ell}}(1)\handwritten{\arrow}[r,"{\text{inertie en }D_i}"]\&[4.2cm]G\handwritten{\arrow}[r,"{\handwritten{\phi}}"]\&H
\end{tikzcd}\(.\)}

Ces classes sont conjuguées entre elles.
